%! Characteristics of Quadratic Functions — Unit 2, Lesson 2.0 — Warm-Up (Answer Key)
% Mirrors warmup/main.tex EXACTLY (12pt, one page): -key in place of -boxes, and every \blank{}
% becomes an \ans{}. Every figure, minipage width and vertical skip is byte-identical, so the two
% files paginate alike. Item 3's last answer is the LESSON'S TARGET MISCONCEPTION — debrief it in
% one sentence and leave it hanging; the hook's vote is thirty seconds away.
\documentclass[12pt]{article}
\usepackage{algebra2-article}
\usepackage{algebra2-key}   % pulls in -boxes; do NOT also load -boxes
% Coordinate grid: {xmin}{xmax}{ymin}{ymax}
\newcommand{\lgrid}[4]{%
  \draw[step=1, linegray!40, very thin] (#1,#3) grid (#2,#4);
  \draw[->, semithick, charcoal] (#1,0)--(#2,0) node[below right, font=\scriptsize]{$x$};
  \draw[->, semithick, charcoal] (0,#3)--(0,#4) node[left, font=\scriptsize]{$y$};}
\newcommand{\fdot}[3]{\fill[#1] (#2,#3) circle (3.2pt);}

\begin{document}

\pageheader{Unit 2, Lesson 2.0}{Warm-Up --- Answer Key}
% NAMESTRIP: no \namedateperiod here — mirrors the blank.

\vspace{0.10in}

\boxguard[12]
\begin{notesbox}{Getting Ready: Skills We'll Use Today}
Three things you already know. Today we meet a graph that \emph{turns around}, and these are the
tools you'll need --- work quickly, then check with a neighbour.

\vspace{6pt}
\textbf{1. Evaluate the rule.} For $f(x)=x^2-2x-3$, find each output. Watch the negatives.
\begin{center}
\renewcommand{\arraystretch}{1.3}
\begin{tabular}{|c|c|c|c|c|}
\hline
$x$ & $-1$ & $0$ & $1$ & $3$ \\ \hline
$f(x)=x^2-2x-3$ & \ans{$0$} & \ans{$-3$} & \ans{$-4$} & \ans{$0$} \\ \hline
\end{tabular}
\end{center}
Which two inputs gave an output of $0$? \ans{$-1$ and $3$}

\vspace{6pt}
\textbf{2. Read the line (Lesson 1.1).} The graph of $y=-2x+4$ is drawn for you.
\begin{minipage}[c]{0.60\linewidth}
  $x$-intercept: \ans{$(2,0)$} \qquad $y$-intercept: \ans{$(0,4)$} \\[5pt]
  Is the line increasing or decreasing? \ans{decreasing} \\[5pt]
  Domain: \ans{all real numbers}
\end{minipage}\hfill
\begin{minipage}[c]{0.36\linewidth}\centering
\begin{tikzpicture}[scale=0.40]
  \lgrid{-1.5}{3.5}{-2.5}{5.5}
  \draw[very thick, forest] (-0.5,5)--(3,-2);
  \fdot{forest}{2}{0}
  \fdot{forest}{0}{4}
\end{tikzpicture}
\end{minipage}

\vspace{6pt}
\textbf{3. A graph that turns around.} This one is \emph{not} a line. Its highest point is
circled.
\begin{minipage}[c]{0.60\linewidth}
  It is increasing until $x=$ \ans{$2$}, and decreasing after $x=$ \ans{$2$}. \\[5pt]
  Its highest point is at \ans{$(2,4)$} \\[5pt]
  \textbf{Is the ``highest'' answer the $x$-value, the $y$-value, or both?} \\[3pt]
  \ans{the $y$-value, $4$ --- $2$ says \emph{where}. (Leave it hanging.)}
\end{minipage}\hfill
\begin{minipage}[c]{0.36\linewidth}\centering
\begin{tikzpicture}[scale=0.40]
  \lgrid{-1.5}{5.5}{-2.5}{5.5}
  \draw[very thick, forest, domain=-0.4:4.4, samples=60, smooth]
    plot (\x, {-(\x-2)*(\x-2)+4});
  \fdot{forest}{2}{4}
  \draw[very thick, redacc] (2,4) circle (7pt);
\end{tikzpicture}
\end{minipage}
\end{notesbox}

\end{document}
